% Compile with: uplatex file.tex && dvipdfmx file.dvi
%           or: latexmk -uplatex file.tex
%\documentclass[uplatex,dvipdfmx]{jsarticle}
\documentclass[dvipdfmx,a4paper,11pt]{jsarticle}
\usepackage{amsmath, amssymb}
\usepackage{tikz}

\title{scenario\_hyperfine\_structure}
\date{\today}

\begin{document}
\maketitle

超微細相互作用 $H_{\rm hf}=A\,\mathbf I\cdot\mathbf S$ を2つの独立な経路で扱う。$A$ は超微細結合定数。

A をスカラーパラメータとして定義して

\[
  A
\]

〔経路A〕核スピン $I=3/2$ と電子スピン $S=1/2$ を合成する。許される全角運動量 $F$ と、$\mathbf I\cdot\mathbf S$ を $F$ で表す恒等式が得られる。$F$ は記号のまま残るので、このあと通常の一般操作に乗せられる。

I=3/2 と S=1/2 を合成して

\begin{equation}
  \left| \frac{3}{2} \right\rangle \otimes \left| \frac{1}{2} \right\rangle:\quad F = 1,\; 2
  \label{eq:E1}
\end{equation}

\begin{equation}
  \mathbf{I}\cdot\mathbf{S} = \frac{1}{2}\left[ F(F+1) - \frac{15}{4} - \frac{3}{4} \right]
  \label{eq:E2}
\end{equation}

恒等式の両辺に結合定数 $A$ を掛けて、超微細ハミルトニアンの固有値の形にする。CAS を経由せずに作った恒等式が、通常の両辺操作にそのまま乗ることの確認でもある。

E2 の両辺に A をかける

\begin{equation}
  A \mathbf{I}\cdot\mathbf{S} = \frac{2 A F^2+2 A F-9 A}{4}
  \label{eq:E3}
\end{equation}

上の準位 $F=2$ を代入する。

E3 の F に 2 を代入する

\begin{equation}
  A \mathbf{I}\cdot\mathbf{S} = \frac{3 A}{4}
  \label{eq:E4}
\end{equation}

下の準位 $F=1$ を代入する。

E3 の F に 1 を代入する

\begin{equation}
  A \mathbf{I}\cdot\mathbf{S} = -\frac{5 A}{4}
  \label{eq:E5}
\end{equation}

2準位の間隔を取る。ランデの間隔則 $\Delta E = A\,F_{\rm max}$ より $2A$ になるはず。

$E4 - E5$

\begin{equation}
  E4 - E5
  \label{eq:E6}
\end{equation}

E6 を簡約して

\begin{equation}
  2 A
  \label{eq:E7}
\end{equation}

〔経路B〕今度はスピン代数のトレースで、同じ $\mathbf I\cdot\mathbf S$ の無限温度モーメントを求める。まず各サイトのスピンの大きさを宣言する——どのスピンでトレースを取ったかは結果の意味の一部なので、これはタイムラインのステップとして残る。

サイト1のスピンを1/2、サイト2のスピンを3/2とする

\noindent\textbf{Assumption:}\ $\displaystyle $\par

$\mathbf I\cdot\mathbf S$ を $x,y,z$ 成分の和として登録する。

$S^x_1 S^x_2 + S^y_1 S^y_2 + S^z_1 S^z_2$

\begin{equation}
  S^x_1 S^x_2 + S^y_1 S^y_2 + S^z_1 S^z_2
  \label{eq:E8}
\end{equation}

1次モーメント（期待値）。無限温度では 0 になるはず。

E8 の期待値

\begin{equation}
  \langle E8 \rangle = 0
  \label{eq:E9}
\end{equation}

2次モーメント。$\frac13 x_1x_2$（$x_i=S_i(S_i+1)$）に $x_1=3/4,\ x_2=15/4$ を入れて $15/16$ になるはず——サイトごとに違うスピンが効いていることの確認。

E8 の2次のモーメント

\begin{equation}
  \langle {E8}^{2} \rangle = \frac{15}{16}
  \label{eq:E10}
\end{equation}

\end{document}